% !TeX root = Tesis.tex

\begin{resumen}
Uno de los desafíos en las áreas de computación e ingeniería es el desarrollo de sistemas que permitan la automatización de los procesos que se generan durante la ejecución de pruebas experimentales que requieren información bioeléctrica (estudios sobre emociones, memoria cognitiva, demencia, neuromarketing, mencionar algunos de ellos). Además, representan un problema a los usuarios que utilizan de este tipo de sistemas, ya que generalmente para adquirirlos se requiere del pago de una licencia y en consecuencia generan dependencia con el proveedor. El objetivo de este trabajo de investigación es diseñar, desarrollar y evaluar un sistema modular empleando software libre que permita obtener, almacenar, administrar y visualizar la información bioeléctrica obtenida de los sensores en tiempo real (un electroencefalógrafo, un sensor de ritmo cardíaco y un sensor de conductancia de la piel), información del expediente clínico electrónico, así como información adicional generada durante la ejecución de un experimento. Para la autenticación del sistema se cuenta con tres tipos de usuarios con privilegios específicos y dos niveles de seguridad, uno por medio de usuario y contraseña, y el otro por medio de una huella digital. Por otra parte, el sistema cuenta con la opción de exportar la información bioeléctrica de los experimentos en formato csv y reportes generales como el expediente clínico electrónico en formato pdf. El sistema que se presenta fue desarrollado empleando software libre basado en el modelo de procesos de ingeniería de la usabilidad y la accesibilidad para el diseño de interfaces, el modelo de proceso evolutivo para la mejora continua de prototipos y el Modelo-Vista-Controlador en la implementación de interfaces gráficas de usuario. Para la evaluación del sistema se realizaron dos tipos de pruebas, las pruebas de funcionalidad para validar los requerimientos y las pruebas de evaluación heurística para evaluar el nivel de usabilidad que presenta el sistema. Finalmente, se concluye que al desarrollar un sistema integrando diferentes tecnologías, empleando software libre y siguiendo un conjunto de metodologías, es posible obtener un sistema funcional y usable, para la automatización de la información en la ejecución de pruebas experimentales que requieren de información bioeléctrica.
%Uno de los desafíos en las áreas de la informática e ingeniería es el desarrollo de sistemas que permitan la automatización de los procesos que se generan durante la ejecución de pruebas experimentales que requieran de información bioeléctrica (estudios sobre las emociones, memoria cognitiva, demencia, neuromarketing, por mencionar algunas).
%Sistemas que permitan la automatización de los procesos que se generan durante la ejecución de pruebas experimentales (estudios sobre las emociones, memoria cognitiva, demencia, neuromarketing, entre otras) que requieren de información bioeléctrica de un individuo, representan un desafío en las áreas de informática e ingeniería principalmente, ya que se precisa de la integración y de la configuración entre el software (bases de datos, drivers, redes de comunicación, aplicaciones, entre otros) y el hardware (sensores, dispositivos de entrada y salida de información, dispositivos de interacción, entre otros). Por otra parte, representan un problema a los usuarios que utilizan de este tipo de sistemas, ya que generalmente para adquirirlos se requiere del pago de una licencia y en consecuencia generan dependencia con el proveedor. El objetivo de este trabajo de investigación es diseñar, desarrollar y evaluar un sistema modular empleando software libre que permita obtener, almacenar, administrar y visualizar la información bioeléctrica obtenida de los sensores en tiempo real (un electroencefalógrafo, un sensor de ritmo cardíaco y un sensor de conductancia de la piel), información del expediente clínico electrónico, así como información adicional generada durante la ejecución de un experimento. Para la autenticación del sistema se cuenta con tres tipos de usuarios con privilegios específicos y dos niveles de seguridad, uno por medio de usuario y contraseña, y el otro por medio de una huella digital. Por otra parte, el sistema cuenta con la opción de exportar la información bioeléctrica de los experimentos en formato csv y reportes generales como el expediente clínico electrónico en formato pdf. El sistema que se presenta fue desarrollado empleando software libre basado en el modelo de procesos de ingeniería de la usabilidad y la accesibilidad para el diseño de interfaces, el modelo de proceso evolutivo para la mejora continua de prototipos y el Modelo-Vista-Controlador en la implementación de interfaces gráficas de usuario. Para la evaluación del sistema se realizaron dos tipos de pruebas, las pruebas de funcionalidad para validar los requerimientos y las pruebas de evaluación heurística para evaluar el nivel de usabilidad que presenta el sistema. Finalmente, se concluye que al desarrollar un sistema integrando diferentes tecnologías, empleando software libre y siguiendo un conjunto de metodologías, es posible obtener un sistema funcional y usable, para la automatización de la información en la ejecución de pruebas experimentales.
%Uno de los desafíos en las áreas de computación e ingeniería es el desarrollo de sistemas que pueden modificar la automatización de los procesos que se pueden realizar durante la ejecución de pruebas experimentales que requieren información bioeléctrica (estudios sobre emociones, memoria cognitiva, demencia, neuromarketing, mencionar algunos de ellos). Además, para los usuarios que usan este tipo de sistema, representa un problema, ya que en general, para comprarlos, se requiere una licencia y, en consecuencia, se genera una dependencia del proveedor. El objetivo de este trabajo de investigación es diseñar, desarrollar y evaluar un sistema modular utilizando un software gratuito que permita obtener, especificar, gestionar y visualizar la información bioeléctrica obteniendo los sensores en tiempo real (un electroencefalograma, un sensor de frecuencia cardíaca y una conductancia de la piel) , información del registro clínico electrónico, así como información adicional generada durante la ejecución de un experimento. Para la autenticación del sistema, hay tres tipos de usuarios con privilegios específicos y dos niveles de seguridad, el primero mediante el uso de un nombre de usuario y contraseña, y el otro mediante huella digital. El sistema también tiene la opción de exportar la información bioeléctrica de los experimentos en formato csv e informes generales, como el registro clínico electrónico en formato pdf. El sistema presentado se desarrolló utilizando software gratuito basado en el modelo de proceso de ingeniería de usabilidad y accesibilidad para el diseño de interfaces, el modelo de proceso evolutivo para la mejora continua de prototipos y el Modelo-Vista-Controlador en la implementación de interfaces gráficas de usuario. Para la evaluación del sistema, se utilizan dos tipos de pruebas: las pruebas de funcionalidad para validar los requisitos y las pruebas de evaluación heurísticas para evaluar el nivel de usabilidad que presenta el sistema. Finalmente, se concluye que al desarrollar un sistema que integra diferentes tecnologías, utilizando software libre y siguiendo un conjunto de metodologías, es posible obtener un sistema funcional y utilizable para la automatización de la información en la ejecución de pruebas experimentales que requieren información bioeléctrica.
\end{resumen}
\begin{abstract}
One of the challenges in the areas of computing and engineering is the development of systems that can modify the automation of the processes that can be performed during the execution of experimental tests that require bioelectrical information (studies on emotions, cognitive memory, dementia, neuromarketing, to mention some of them). Morover, for users who use this type-kind of system, it represents a problem, since in general, to purchase them, a license is required and, consequently, it is generated a dependence on the provider. The objective of this research work is to design, develop and evaluate a modular system using free software that allows obtaining, specifying, managing and visualizing the bioelectrical information obtaining the sensors in real time (an electroencephalograph, a heart rate sensor and a skin conductance), information from the electronic clinical record, as well as additional information generated during the execution of an experiment. For system authentication, there are three types of users with specific privileges and two levels of security, the first by using an username and password, and the other via fingerprint. The system has too the option of exporting the bioelectrical information of the experiments in csv format and general reports such as the electronic clinical record in pdf format. The system presented was developed using free software based on the usability and accessibility engineering process model for interface design, the evolutionary process model for continuous improvement of prototypes and the Model-View-Controller in the implementation of graphical user interfaces. For the evaluation of the system, two types of tests are used: the functionality tests to validate the requirements and the heuristic evaluation tests to evaluate the level of usability that the system presents. Finally, it is concluded that by developing a system integrating different technologies, using free software and following a set of methodologies, it is possible to obtain a functional and usable system for the automation of information in the execution of experimental tests that requiers biolectrical information.
%Systems that allowed the automation of the processes that were carried out during the execution of experimental tests (studies on emotions, cognitive memory, dementia, neuromarketing, among others) and engineering mainly, since it requires integration and configuration between software (databases, controllers, communication networks, applications, among others) and hardware (sensors, information input and output devices, interaction devices, among others). On the other hand, it represents a problem for users who use this type of system, since in general, to purchase them, a license is required and, consequently, dependence on the provider depends. The objective of this research work is to design, develop and evaluate a modular system using free software that allows obtaining, specifying, managing and visualizing the bioelectrical information obtaining the sensors in real time (an electroencephalograph, a heart rate sensor and a skin conductance), information from the electronic clinical record, as well as additional information generated during the execution of an experiment. For system authentication, there are three types of users with specific privileges and two levels of security, one by means of username and password, and the other by means of a fingerprint. On the other hand, the system has the option of exporting the bioelectrical information of the experiments in csv format and general reports such as the electronic clinical record in pdf format. The system presented was developed using free software based on the usability and accessibility engineering process model for interface design, the evolutionary process model for continuous improvement of prototypes and the Model-View-Controller in the implementation of graphical user interfaces. Two types of tests are used for the evaluation of the system, the functionality tests to validate the requirements and the heuristic evaluation tests to evaluate the level of usability that the system presents. Finally, it is concluded that by developing a system integrating different technologies, using free software and following a set of methodologies, it is possible to obtain a functional and usable system for the automation of information in the execution of experimental tests.


%Systems that allow the automation of the processes carried out during the execution of experimental tests (studies on emotions, cognitive memory, dementia, neuromarketing, among others) that require bioelectrical information from an individual, represent a challenge in the areas of computing and engineering mainly, since it needs the integration and configuration between the software (databases, controllers, communication networks, applications, among others) and the hardware (sensors, information input and output devices, interaction devices, among others). On the other hand, represent a problem to users who need this type of systems, since generally to acquire them requires the payment of a license and consequently depend on dependence on the provider. Consequently, the objective of this research work is to develop, develop and evaluate a modular system using free software that allows obtaining, controlling, managing and visualizing the bioelectrical information from the sensors in real time (an electroencephalograph, a heart rate sensor and a skin conductance sensor), information from the electronic clinical record, as well as additional information generated during the execution of an experiment. For system authentication, there are three types of users with specific privileges and two levels of security, one by means of username and password, and the other by means of a fingerprint. On the other hand, the system has the option of exporting the bioelectrical information of the experiments in csv format and general reports such as the electronic clinical record in pdf format. This system was developed using free software based on the usability and accessibility engineering process model for interface design, the evolutionary process model for continuous prototype improvement, and the controller view model in the implementation of graphical interfaces of user. Two types of tests are used for the evaluation of the system, the functionality tests to validate the system requirements and the heuristic evaluation tests to evaluate the level of usability that the system presents. Finally, it is concluded that by generating a system integrating different technologies, using free software and following a set of methodologies, it is possible to obtain a functional and usable system for the automation of information in the execution of experimental tests.
%The automation of the processes that are carried out during the experimental tests that require the bioelectric information of an individual represents a challenge in the areas of computer science and engineering mainly, since it needs the integration and the configuration between the software (databases, controllers, communication networks, applications, among others) and hardware (sensors, information input and output devices, interaction devices, among others). In the market there are systems that have the objective but, they represent a problem when acquiring it for its high cost and also a dependency with the supplier is generated. The Bioelectric Information Based Information System (SIBIB) is a modular system, where the part of it allows obtaining, specifying and visualizing the bioelectric information of three sensors in real time produced by users during the interaction with a product or multimedia media. These are an electroencephalograph, a heart rate sensor and an electrical skin conductance sensor. Also, the system has two levels of security for access authentication, one by means of a username and password, and the other by means of a fingerprint. In comparison with other systems that carry out this process, it also allows the visualization of the information collected by the sensors in real time, its storage, the addition of comments and markers. There is also the option to export the information of the experiments and general reports such as the electronic clinical file. This system was developed using free software based on the usability and accessibility engineering process model for the design of its interfaces, the evolutionary process model for continuous prototype improvement and the controller view model for code implementation. of programming. For its evaluation two types of tests are used, the functionality tests and a heuristic evaluation, which is carried out by three evaluators. Finally, it is concluded that by generating a system integrating different technologies, a functional and usable system can be obtained for experimental tests that require bioelectric information.
%The realization of experimental tests to obtain bioelectrical information of an individual is a challenge for different investigations or projects, this is due to two factors that are necessary for the realization of these. The first one is the software and the second the hardware. In most cases you cannot use one without having access to the other. For this reason there are companies that focus on developing both the hardware for testing and the software for its use. The Information System Based on Bioelectric Information is a set of various technologies that deal with the collection, storage and visualization of real-time information produced by users during interaction with a product or multimedia media. The system consists of the integration of three sensors for the collection of bioelectric information. These are an electroencephalograph, a heart rate sensor and an electrical skin conductance sensor. The system also has two levels of authentication for access, one through a username and password, and the other through a fingerprint. Compared to other softwares that perform this process, this allows the visualization of the information collected by the sensors in real time, its storage, the addition of comments and bookmarks in files. There is also the option to export the information of the experiments and general reports such as the electronic clinical file. This system was developed based on the usability and accessibility engineering process model (MPIu + a).
%The Information System Based on Biometric Information is a set of various technologies that deal with the collection, storage and visualization of real-time information produced by biometric experiments. The system consists of the integration of three sensors for the collection of the bioelectrical information that a person produces during the interaction with some element, be it a product, a software or some multimedia element. The sensors integrated in the system are an electroencephalograph, a heart rate sensor and an electrical skin conductance sensor. The system has two levels of authentication for access, one through a username and password, and the other through a fingerprint. It allows the visualization of the information collected by the sensors in real time, its storage and the addition of comments and bookmarks in files. Also, there is the option to export the information of the experiments and general reports. This system was developed based on the usability and accessibility engineering process model.

\end{abstract}

\newpage
\thispagestyle{empty}

%\begin{publicaciones}

%\noindent Como parte de los resultados del trabajo de investigación
%desarrollado en esta tesis, se obtuvo el siguiente artículo (ejemplos para artículos resultado de tesis)s:

%\subsection*{Artículos en congresos internacionales}
%\subsubsection*{Aceptados}
%\begin{itemize}
%\item \textbf{González, L. A., Jarillo, A., Cruz, J.A., Gómez, V., García, A. M., y Domínguez O. A. (2012).}, \emph{Cartesian control
%application in haptic interfaces for motor rehabilitation purposes.}, IEEE ninth electronics (CERMA), Cuernavaca, Morelos, México.
%\end{itemize}
%\subsubsection*{Sometidos}
%\begin{itemize}
%\item \textbf{González, L. A., Jarillo, A., Cruz, J.A., Gómez, V., García, A. M., y Domínguez O. A. (2013).}, \emph{REyEMO a haptic system for rehabilitation purpuses.}, International Journal of Digital content Technology and its Applications (JDCTA), Korea.
%\end{itemize}

%\end{publicaciones}
